\documentclass[11pt,a4paper]{article}
% Shared preamble for the qi26_25 weekly deliverable reports.
%
% Each weekly report is a short standalone document that inputs this file, so
% that formatting stays consistent across all six without duplicating the
% package list in every one. Change something here and every report follows.
%
% Usage in a weekly report:
%     \documentclass[11pt,a4paper]{article}
%     % Shared preamble for the qi26_25 weekly deliverable reports.
%
% Each weekly report is a short standalone document that inputs this file, so
% that formatting stays consistent across all six without duplicating the
% package list in every one. Change something here and every report follows.
%
% Usage in a weekly report:
%     \documentclass[11pt,a4paper]{article}
%     % Shared preamble for the qi26_25 weekly deliverable reports.
%
% Each weekly report is a short standalone document that inputs this file, so
% that formatting stays consistent across all six without duplicating the
% package list in every one. Change something here and every report follows.
%
% Usage in a weekly report:
%     \documentclass[11pt,a4paper]{article}
%     \input{weekly_preamble}
%     \weekhead{1}{Vision Foundation and Data Encoding}{Week of ...}
%     ... body ...
%     \end{document}

\usepackage[utf8]{inputenc}
\usepackage[T1]{fontenc}
\usepackage[margin=2.5cm]{geometry}
\usepackage{amsmath,amssymb}
\usepackage{booktabs}
\usepackage{graphicx}
\usepackage{enumitem}
\usepackage{listings}
\usepackage{xcolor}
\usepackage[hidelinks]{hyperref}
\usepackage{titlesec}
\usepackage{fancyhdr}

\definecolor{codebg}{RGB}{247,247,250}
\definecolor{codekw}{RGB}{20,90,160}
\definecolor{codecm}{RGB}{110,120,130}
\definecolor{accent}{RGB}{20,90,160}

\lstset{
  backgroundcolor=\color{codebg},
  basicstyle=\ttfamily\footnotesize,
  keywordstyle=\color{codekw}\bfseries,
  commentstyle=\color{codecm}\itshape,
  breaklines=true,
  frame=single,
  rulecolor=\color{gray!30},
  language=Python,
  showstringspaces=false,
  xleftmargin=6pt,
  framexleftmargin=6pt
}

\setlist{itemsep=2pt,topsep=4pt}
\titleformat{\section}{\normalfont\large\bfseries\color{accent}}{\thesection}{0.8em}{}
\titleformat{\subsection}{\normalfont\normalsize\bfseries}{\thesubsection}{0.8em}{}

\pagestyle{fancy}
\fancyhf{}
\fancyfoot[C]{\small\thepage}
\fancyfoot[R]{\scriptsize\texttt{github.com/Pushkar0997/qi26\_25}}
\renewcommand{\headrulewidth}{0pt}

% \weekhead{number}{title}{deliverable as stated in the project brief}
\newcommand{\weekhead}[3]{%
  \begin{center}
    {\scriptsize\textsf{QINTERN 2026 \textbullet\ QWORLD \textbullet\ WEEKLY DELIVERABLE}}\\[0.25cm]
    {\LARGE\bfseries Week #1: #2}\\[0.35cm]
    {\normalsize Hybrid QML and Quantum String Algorithms for Document Extraction}\\[0.25cm]
    {\small Pushkar Kumar \quad\textbullet\quad Mentor: Potluri Krishna Priyatham}\\[0.4cm]
    \fbox{\begin{minipage}{0.88\textwidth}
      \small\textbf{Deliverable as specified in the project brief:} #3
    \end{minipage}}
  \end{center}
  \vspace{0.4cm}
}

% A boxed note for caveats and honest limitations.
\newcommand{\honestnote}[1]{%
  \vspace{0.2cm}
  \noindent\fcolorbox{gray!40}{gray!8}{%
    \begin{minipage}{0.97\textwidth}\small #1\end{minipage}}
  \vspace{0.2cm}
}

%     \weekhead{1}{Vision Foundation and Data Encoding}{Week of ...}
%     ... body ...
%     \end{document}

\usepackage[utf8]{inputenc}
\usepackage[T1]{fontenc}
\usepackage[margin=2.5cm]{geometry}
\usepackage{amsmath,amssymb}
\usepackage{booktabs}
\usepackage{graphicx}
\usepackage{enumitem}
\usepackage{listings}
\usepackage{xcolor}
\usepackage[hidelinks]{hyperref}
\usepackage{titlesec}
\usepackage{fancyhdr}

\definecolor{codebg}{RGB}{247,247,250}
\definecolor{codekw}{RGB}{20,90,160}
\definecolor{codecm}{RGB}{110,120,130}
\definecolor{accent}{RGB}{20,90,160}

\lstset{
  backgroundcolor=\color{codebg},
  basicstyle=\ttfamily\footnotesize,
  keywordstyle=\color{codekw}\bfseries,
  commentstyle=\color{codecm}\itshape,
  breaklines=true,
  frame=single,
  rulecolor=\color{gray!30},
  language=Python,
  showstringspaces=false,
  xleftmargin=6pt,
  framexleftmargin=6pt
}

\setlist{itemsep=2pt,topsep=4pt}
\titleformat{\section}{\normalfont\large\bfseries\color{accent}}{\thesection}{0.8em}{}
\titleformat{\subsection}{\normalfont\normalsize\bfseries}{\thesubsection}{0.8em}{}

\pagestyle{fancy}
\fancyhf{}
\fancyfoot[C]{\small\thepage}
\fancyfoot[R]{\scriptsize\texttt{github.com/Pushkar0997/qi26\_25}}
\renewcommand{\headrulewidth}{0pt}

% \weekhead{number}{title}{deliverable as stated in the project brief}
\newcommand{\weekhead}[3]{%
  \begin{center}
    {\scriptsize\textsf{QINTERN 2026 \textbullet\ QWORLD \textbullet\ WEEKLY DELIVERABLE}}\\[0.25cm]
    {\LARGE\bfseries Week #1: #2}\\[0.35cm]
    {\normalsize Hybrid QML and Quantum String Algorithms for Document Extraction}\\[0.25cm]
    {\small Pushkar Kumar \quad\textbullet\quad Mentor: Potluri Krishna Priyatham}\\[0.4cm]
    \fbox{\begin{minipage}{0.88\textwidth}
      \small\textbf{Deliverable as specified in the project brief:} #3
    \end{minipage}}
  \end{center}
  \vspace{0.4cm}
}

% A boxed note for caveats and honest limitations.
\newcommand{\honestnote}[1]{%
  \vspace{0.2cm}
  \noindent\fcolorbox{gray!40}{gray!8}{%
    \begin{minipage}{0.97\textwidth}\small #1\end{minipage}}
  \vspace{0.2cm}
}

%     \weekhead{1}{Vision Foundation and Data Encoding}{Week of ...}
%     ... body ...
%     \end{document}

\usepackage[utf8]{inputenc}
\usepackage[T1]{fontenc}
\usepackage[margin=2.5cm]{geometry}
\usepackage{amsmath,amssymb}
\usepackage{booktabs}
\usepackage{graphicx}
\usepackage{enumitem}
\usepackage{listings}
\usepackage{xcolor}
\usepackage[hidelinks]{hyperref}
\usepackage{titlesec}
\usepackage{fancyhdr}

\definecolor{codebg}{RGB}{247,247,250}
\definecolor{codekw}{RGB}{20,90,160}
\definecolor{codecm}{RGB}{110,120,130}
\definecolor{accent}{RGB}{20,90,160}

\lstset{
  backgroundcolor=\color{codebg},
  basicstyle=\ttfamily\footnotesize,
  keywordstyle=\color{codekw}\bfseries,
  commentstyle=\color{codecm}\itshape,
  breaklines=true,
  frame=single,
  rulecolor=\color{gray!30},
  language=Python,
  showstringspaces=false,
  xleftmargin=6pt,
  framexleftmargin=6pt
}

\setlist{itemsep=2pt,topsep=4pt}
\titleformat{\section}{\normalfont\large\bfseries\color{accent}}{\thesection}{0.8em}{}
\titleformat{\subsection}{\normalfont\normalsize\bfseries}{\thesubsection}{0.8em}{}

\pagestyle{fancy}
\fancyhf{}
\fancyfoot[C]{\small\thepage}
\fancyfoot[R]{\scriptsize\texttt{github.com/Pushkar0997/qi26\_25}}
\renewcommand{\headrulewidth}{0pt}

% \weekhead{number}{title}{deliverable as stated in the project brief}
\newcommand{\weekhead}[3]{%
  \begin{center}
    {\scriptsize\textsf{QINTERN 2026 \textbullet\ QWORLD \textbullet\ WEEKLY DELIVERABLE}}\\[0.25cm]
    {\LARGE\bfseries Week #1: #2}\\[0.35cm]
    {\normalsize Hybrid QML and Quantum String Algorithms for Document Extraction}\\[0.25cm]
    {\small Pushkar Kumar \quad\textbullet\quad Mentor: Potluri Krishna Priyatham}\\[0.4cm]
    \fbox{\begin{minipage}{0.88\textwidth}
      \small\textbf{Deliverable as specified in the project brief:} #3
    \end{minipage}}
  \end{center}
  \vspace{0.4cm}
}

% A boxed note for caveats and honest limitations.
\newcommand{\honestnote}[1]{%
  \vspace{0.2cm}
  \noindent\fcolorbox{gray!40}{gray!8}{%
    \begin{minipage}{0.97\textwidth}\small #1\end{minipage}}
  \vspace{0.2cm}
}


\begin{document}
\weekhead{6}{Benchmarking and Final Reporting}
{Quantitative comparison against classical architectures; final report and presentation.}

\section{Objective}

Compare compute resources against classical approaches, evaluate accuracy against
model size, and produce the final report.

\section{Experimental Design}

Five experiments were defined, all reproducible from a committed dataset:

\begin{enumerate}[label=\Alph*.]
\item Feature-extractor comparison, per tier, at matched dimensionality.
\item End-to-end character error rate and identifier recovery per tier.
\item Quantum resource cost: encoding, quanvolutional layer, Grover oracle.
\item Grover oracle scaling with text length.
\item Shot-noise sensitivity of the quantum feature stage.
\end{enumerate}

One design rule governs all of them: \textbf{when comparing quantum against
classical, hold feature dimensionality fixed}. Otherwise a win is
indistinguishable from granting one method a wider representation.

\section{Experiment A: Feature Extractors}

\begin{table}[h]
\centering
\caption{Character accuracy by extractor and tier.}
\small
\begin{tabular}{@{}lrrrr@{}}
\toprule
Tier & quanv (64) & quanv$+ZZ$ (160) & classical (160) & raw pixels (64) \\
\midrule
clean\_digital & 99.3\% & 100\% & 100\% & 100\% \\
clean\_scan & 99.3\% & 100\% & 100\% & 100\% \\
noisy\_scan & 93.0\% & 97.7\% & 97.7\% & 98.3\% \\
clear\_handwriting & 87.7\% & 96.7\% & 96.3\% & 98.7\% \\
degraded\_handwriting & 62.5\% & 75.7\% & 74.4\% & 79.4\% \\
\midrule
\textbf{All tiers} & \textbf{86.7\%} & \textbf{92.2\%} & \textbf{91.9\%} & \textbf{93.6\%} \\
\bottomrule
\end{tabular}
\end{table}

\subsection{Significance}

A single split gives point estimates but cannot separate a real difference from
noise. The comparison was repeated across 30 splits, with every extractor scored
on the same splits so comparisons are paired.

\begin{table}[h]
\centering
\caption{Paired comparisons across 30 splits.}
\small
\begin{tabular}{@{}lrrl@{}}
\toprule
Comparison & Difference & $p$ & Verdict \\
\midrule
quanv$+ZZ$ $-$ classical & $-0.1$ pt & $\approx 0.2$--$0.3$ & not distinguishable from noise \\
raw pixels $-$ quanv$+ZZ$ & $+1.2$ pt & $< 10^{-10}$ & significant \\
raw pixels $-$ classical & $+1.1$ pt & $< 10^{-12}$ & significant \\
quanv$+ZZ$ $-$ quanv & $+5.7$ pt & $< 10^{-25}$ & significant \\
\bottomrule
\end{tabular}
\end{table}

\honestnote{%
Figures are given as ranges because scikit-learn's solver is not bit-reproducible
across BLAS builds: repeated runs on different machines give $p$ between roughly
0.2 and 0.3 and win counts of 14 or 15 out of 30. The conclusions are stable
across every run; the third decimal place is not, and quoting it would imply a
precision this measurement does not have.}

This separates two claims that the single-split table presents identically. The
quantum-versus-classical gap is noise---the quantum layer wins roughly half the
splits. The raw-pixel lead is not: it holds in all 30.

\section{Experiment B: End-to-End}

\begin{table}[h]
\centering
\caption{End-to-end performance, all 12 documents per tier.}
\begin{tabular}{@{}lrrr@{}}
\toprule
Tier & CER & Segmentation ratio & ID recovered exactly \\
\midrule
clean\_digital & 6.6\% & 0.95 & \textbf{100\%} \\
clean\_scan & 8.5\% & 0.93 & 58\% \\
noisy\_scan & 79.5\% & 0.34 & 0\% \\
clear\_handwriting & 51.8\% & 0.70 & 0\% \\
degraded\_handwriting & 85.7\% & 0.99 & 0\% \\
\bottomrule
\end{tabular}
\end{table}

The identifier column is the one that matters. Character error rate measures
average quality; identifier recovery measures whether the pipeline did its job,
for which one wrong character makes the result useless. The two diverge sharply
on \texttt{clean\_scan}: 8.5\% CER but only 58\% usable identifiers.

\section{Experiments C and D: Resource Cost and Scaling}

\begin{table}[h]
\centering
\caption{Quantum resource cost.}
\begin{tabular}{@{}lrrr@{}}
\toprule
Circuit & Qubits & Depth & CX gates \\
\midrule
Quanvolutional filter (per patch) & 4 & 10 & 8 \\
Grover oracle, 16-character field & 14 & 6{,}073 & 3{,}272 \\
\bottomrule
\end{tabular}
\end{table}

The filter runs 16 times per character, so the per-character cost is 128 CX
gates; the per-document cost scales with character count.

Oracle scaling was measured across text lengths from 8 to 64 characters. Fitting
a power law gives a CX exponent of $\mathbf{1.39}$: superlinear, because each
candidate position requires its own controlled data load. The measured curve
sits above a linear reference and nowhere near $\sqrt{N}$.

\section{Experiment E: Shot Noise}

\begin{table}[h]
\centering
\caption{Accuracy against shot budget per patch.}
\begin{tabular}{@{}lr@{}}
\toprule
Shots & Accuracy \\
\midrule
64 & 70.7\% \\
256 & 81.7\% \\
1{,}024 & 85.2\% \\
4{,}096 & 86.7\% \\
exact & 86.7\% \\
\bottomrule
\end{tabular}
\end{table}

Approximately 1{,}024 shots per patch approach the noiseless result. At 16
patches per character this is roughly 16{,}000 circuit executions per
character---a substantial hardware cost that belongs in any honest resource
accounting and is frequently absent from published quanvolutional results.

\section{External Baseline}

Experiments A and B compare feature maps that all share this project's
segmentation front end and classifier, which answers ``which feature map is best
inside this pipeline?'' but leaves the pipeline unanchored. Tesseract provides
the external reference, run on the same images over the same twelve documents per
tier.

\begin{table}[h]
\centering
\caption{Comparison against Tesseract.}
\begin{tabular}{@{}lrrrr@{}}
\toprule
Tier & This pipeline & Tesseract & Pipeline ID & Tesseract ID \\
\midrule
clean\_digital & 6.6\% & \textbf{2.0\%} & \textbf{100\%} & 92\% \\
clean\_scan & 8.5\% & \textbf{1.9\%} & 58\% & \textbf{83\%} \\
noisy\_scan & \textbf{79.5\%} & 100.0\% & 0\% & 0\% \\
clear\_handwriting & 51.8\% & \textbf{20.2\%} & 0\% & \textbf{17\%} \\
degraded\_handwriting & \textbf{85.7\%} & 100.0\% & 0\% & 0\% \\
\bottomrule
\end{tabular}
\end{table}

Tesseract is three to four times better on clean input. That is the expected
result and is reported plainly. Two details complicate rather than rescue the
picture: on \texttt{clean\_digital} this pipeline recovers more identifiers,
because a constrained 36-class charset cannot emit out-of-charset characters; and
on the two most degraded tiers Tesseract returns nothing at all, its page
analysis rejecting input it judges unreadable.

\honestnote{%
This comparison is not controlled---Tesseract brings its own segmentation,
language model and training corpus. It is included because the first question a
sceptical reader asks is how the pipeline compares to an ordinary classical
tool, and the answer belongs in the report regardless of whether it flatters the
work.}

\section{Accuracy Against Model Size}

The brief asked for accuracy against model size, in the spirit of comparing QNN
parameters against LLM weights.

\begin{table}[h]
\centering
\caption{Parameter count against accuracy.}
\begin{tabular}{@{}lrr@{}}
\toprule
Feature stage & Trainable parameters & Accuracy \\
\midrule
Quanvolutional filter & 8 & 91.7\% \\
Classical convolution, dimension-matched & 40 & 92.0\% \\
Raw pixels & 0 & \textbf{93.6\%} \\
\bottomrule
\end{tabular}
\end{table}

The quantum filter is the most parameter-efficient of the three that have
parameters, producing a 160-dimensional feature space from 8 angles. But the
zero-parameter option outperforms both, which is the finding: at this input size
the parameter-efficiency comparison is won by not extracting features at all.

\section{Conclusions}

\begin{enumerate}
\item At matched dimensionality the quantum feature layer is statistically
      indistinguishable from a classical convolution, and both lose to raw
      pixels. Training the filter does not change this.
\item Grover string matching works correctly, but oracle cost grows as
      $N^{1.39}$ owing to data loading, so the $\sqrt{N}$ query advantage does
      not become an end-to-end speedup for stored classical text.
\item The dominant error source is the classical segmentation front end, not
      either quantum stage.
\item Against Tesseract, this pipeline is three to four times worse on clean
      input. It is not competitive with mature classical OCR.
\item The single largest improvement came from correcting a train/serve data
      skew, not from any modelling change.
\end{enumerate}

The work does not demonstrate quantum advantage for document intelligence, and
identifies specific measured reasons why. Both are useful negative results, and
both are more informative than a tuned demonstration would have been.

\section{Reproducibility}

Every figure regenerates from a clean repository clone. The dataset and trained
model are committed; the browser demonstration is verified byte-identical to the
Python implementation; and continuous integration re-runs the benchmark on every
push and fails the build if any figure in the results file does not appear in the
report. That final gate has caught a real contamination---a results file measured
on a drifted dataset.

\section{Deliverable Status}

Complete. Five experiments implemented and measured, significance established,
external baseline added, final report produced.

\vspace{0.3cm}
\noindent\textbf{Artifacts:} \texttt{benchmarking/run\_benchmark.py},
\texttt{benchmarking/seed\_sweep.py}, \texttt{benchmarking/tesseract\_baseline.py},
\texttt{benchmarking/make\_figures.py}, \texttt{docs/final\_report.md},
\texttt{reports/qi26\_25\_final\_report.pdf}

\end{document}
